%!TEX root = ../main.tex
\newpage
\section{フィードフォワード型三次元再構成}\label{sec:feedforward-3d-reconstruction}

本研究の本筋である最適化ベース(Optimization-based)の手法とは対照的に、近年(2023-2025年)、大規模データセットとDeep Learning、特にTransformerアーキテクチャの発展により、幾何学的計算を推論(Inference)として解く「フィードフォワード型(Feed-Forward)」の手法が台頭している。
従来の手法(SfM/MVSやNeRF/3DGS)が、入力シーンごとにパラメータを反復的に更新して解を探索するのに対し、フィードフォワード型の手法は、学習済みの膨大な事前知識(Priors)を用いて、単一の順伝播処理のみで三次元構造を回帰するデータ駆動型(Data-driven)のアプローチである。

このアプローチの代表例として、以下の手法が挙げられる。

\begin{itemize} 
  \item \textbf{Depth Anything} \cite{Yang2024CVPR_DepthAnything,Yang2024arXiv_DepthAnythingV2}: 
  6200万枚以上の画像から学習された基盤モデル(Foundation Model)であり、DINOv2バックボーンを活用することで、テクスチャのない領域や未知のシーンに対しても極めてロバストな単眼深度推定を実現している。
  しかし、出力はスケール不確定性を伴う2.5次元表現に留まり、多視点間での厳密な幾何学的整合性は保証されない。
  Gaussian Splattingの正則化に用いられる場合もあり、その際はSfM点群によってスケールを持たせたDepth Anythingの深度マップと、Gaussian Splattingの深度レンダリングとの誤差を最小化する。

  \item \textbf{DUSt3R} \cite{Wang2024CVPR_DUSt3R, Leroy2024ECCV_MASt3R}: 
  従来のSfMパイプライン(特徴点抽出、マッチング、バンドル調整)を完全に排除し、2枚の画像から直接「ポイントマップ(Point Map)」を回帰する手法である。 
  これにより、カメラパラメータを事前に与えることなく、エンドツーエンドでの三次元形状およびPoint Mapからのカメラ姿勢の逆推定が可能となった。
  後継のMASt3Rでは、局所的なマッチング精度の向上が図られており、より高精度な三次元アライメントを実現している。

  \item \textbf{VGGT} \cite{Wang2025CVPR_VGGT}: 
  Visual Geometry Grounded Transformer (VGGT)は、任意の枚数の画像を入力とし、カメラパラメータ、深度、点群、そして追跡情報(Tracks)の全てを同時に推論する。
  トランスフォーマーを用いた統合的なアーキテクチャにより、従来の幾何学的制約を学習された高次元の構造的特徴として捉え、ノイズの多い入力や極端な視点変化に対しても極めて高い頑健性を示す。

\end{itemize}

これらの手法は、計算速度と多様な環境における頑健性において革新的であるが、本研究が扱う屈折を含むシーンでは明確な課題もある。
フィードフォワード手法はその性質上、学習データに深く性能を依存する。
一般的な学習データセットには屈折が十分に、かつ物理的に正確なアノテーションと共に含まれていないため、フィードフォワードモデルでは正確な幾何情報を再現できない可能性が考えられる。
対して、本研究で用いるGaussian Splatting等の最適化ベース手法は、屈折率やスネルの法則を明示的にモデル化し組み込むことで、こうした物理現象を正確に逆算することが可能である。
しかし、動画生成AIが流体や光の反射といった物理法則をデータから獲得しつつある現状を鑑みると\cite{Liu2024arXiv_Sora-review}、将来的にはフィードフォワード手法も十分なデータスケールによって屈折を「学習」する可能性は否定できない。

